\section*{Bab \ref{ch:b3:water-rock-interaction} --- Interaksi Air--Batuan}

\begin{solution}{exo:b3:water-rock-interaction:1}
Kation: $2\times48/40.08 + 2\times9.7/24.31 + 23/22.99 + 3.9/39.10 =
\qty{4.29}{meq/L}$; anion: $183/61.02 + 35.5/35.45 + 2\times48/96.06 =
\qty{5.00}{meq/L}$. Galat $100\times(4.29 - 5.00)/9.29 = \qty{-7.6}{\percent}$:
jauh melampaui 1 atau \qty{2}{\percent} yang wajar untuk analisis cermat terhadap air semacam ini.
Satu kation hilang atau kurang dilaporkan, atau satu anion terlalu tinggi.
\end{solution}

\begin{solution}{exo:b3:water-rock-interaction:2}
$\sqrt I = 0.0707$. \ce{Ca^2+}: $-\log\gamma = 0.5085\times4\times0.0707/(1 +
0.3281\times6\times0.0707) = 0.126$, $\gamma = 0.748$; \ce{HCO3-}: $\gamma =
0.927$; \ce{SO4^2-}: $\gamma = 0.739$.
\end{solution}

\begin{solution}{exo:b3:water-rock-interaction:3}
$\sqrt I = 0.126$; $\gamma_{\ce{Ca}} = 0.622$, $\gamma_{\ce{SO4}} = 0.602$;
$\log\mathrm{IAP} = \log(3.0\times10^{-3}\times0.622\times4.0\times10^{-3}
\times0.602) = -5.35$; $\mathrm{SI} = -5.35 + 4.55 = -0.80$. Air ini
belum jenuh dan masih dapat melarutkan gipsum.
\end{solution}

\begin{solution}{exo:b3:water-rock-interaction:4}
$m = 0.012/60.08 = \qty{2.00e-4}{mol/L}$, $\log m = -3.70$. Kuarsa:
$\mathrm{SI} = -3.70 + 3.98 = +0.28$, sedikit lewat jenuh; silika
amorf: $-3.70 + 2.71 = -0.99$, belum jenuh.
\end{solution}

\begin{solution}{exo:b3:water-rock-interaction:5}
Pada \qty{25}{\celsius}, $\log K = -8.45 - 6.35 - 1.47 + 10.33 = -5.94$, $m =
(10^{-5.94}\times10^{-2}/4)^{1/3} = \qty{1.42}{mmol/L}$, \qty{57}{mg/L}; pada
\qty{10}{\celsius}, $\log K = -5.59$, $m = \qty{1.86}{mmol/L}$,
\qty{75}{mg/L}: sekitar sepertiga lebih banyak di air dingin.
\end{solution}

\begin{solution}{exo:b3:water-rock-interaction:6}
$[\ce{CO2}] = 10^{-7.60}\times4.0\times10^{-3}/10^{-6.35} =
\qty{2.25e-4}{mol/L}$; $[\ce{CO3^2-}] = 10^{-10.33}\times4.0\times10^{-3}/
10^{-7.60} = \qty{7.4e-6}{mol/L}$; $\log P_{\ce{CO2}} = -3.648 + 1.47 =
-2.18$, 16 kali lipat dibanding udara. Total \qty{4.23}{mmol/L}: \ce{CO2}
\qty{5.3}{\percent}, \ce{HCO3-} \qty{94.5}{\percent}, \ce{CO3^2-}
\qty{0.18}{\percent}.
\end{solution}

\begin{solution}{exo:b3:water-rock-interaction:7}
$r = 10^{-12.09} = \qty{8.2e-13}{mol.m^{-2}.s^{-1}}$ (suku asam, netral dan
basa masing-masing menyumbang $10^{-12.45}$, $10^{-12.56}$ dan $10^{-12.74}$). Jari-jari
menyusut dengan laju $rV_m = 8.2\times10^{-13}\times1.0131\times10^{-4} =
\qty{8.3e-17}{m/s}$, \qty{2.6}{nm} per tahun; jari-jari \qty{0.5}{mm} bertahan
$6.0\times10^{12}$~s, sekitar \num{190000} tahun.
\end{solution}

\begin{solution}{exo:b3:water-rock-interaction:8}
$\exp[(69.8/0.008314)(1/278.15 - 1/298.15)] = 7.6$; dengan
\qty{74}{kJ/mol}, 8.6.
\end{solution}

\begin{solution}{exo:b3:water-rock-interaction:9}
Dengan menyamakan $(12.703 - X - 3Y)/3$ dan $(7.435 - 2Y)/2$, $Y$ tereliminasi: $X =
12.703 - 11.153 = 1.55$, sebuah garis horizontal. Air tersebut memiliki $X = \log(2.0\times
10^{-3}/39.10) + 7.0 = 2.71$ dan $Y = \log(0.015/60.08) = -3.60$: di atas 1.55
dan jauh di bawah garis muskovit--K-feldspar ($X = -3.161 - 3Y = 7.65$), maka
terbentuk muskovit.
\end{solution}

\begin{solution}{exo:b3:water-rock-interaction:10}
$\delta_w^f = -2 - 2 = -4$. Tertutup: $W/R = 0.5\times(-2 - 6)/(-13 + 4) =
0.44$. Terbuka: $0.5\ln[(-8)/(-9) + 1] = 0.5\ln1.889 = 0.32$. Pada sistem terbuka, setiap
parsel meninggalkan batuan dalam kesetimbangan dengan batuan yang terakhir ditemuinya, sambil membawa
oksigen-18 sebanyak mungkin; dalam sistem tertutup seluruh kumpulan air berakhir pada $-4$ dan
membawa lebih sedikit per kilogram.
\end{solution}

\begin{solution}{exo:b3:water-rock-interaction:11}
Terbuka: $\log K = -5.59$, $m = (10^{-5.59}\times10^{-1.5}/4)^{1/3} =
\qty{2.73}{mmol/L}$, \qty{109}{mg/L}. Tertutup: $C_0 = 10^{-1.27 - 1.5} =
\qty{1.70e-3}{mol/L}$ dan $K_{\mathrm{cal}}K_1/K_2 = 10^{-4.32}$; persamaan kubik
$4m^3 = 4.79\times10^{-5}(1.70\times10^{-3} - m)$ memberikan $m =
\qty{1.45}{mmol/L}$, \qty{58}{mg/L}, sekitar separuhnya. Karbon dioksida yang tersisa
adalah $0.25\times10^{-3}$, \qty{15}{\percent} dari $C_0$, sehingga $\log P_{\ce{CO2}}$ turun
dari $-1.5$ menjadi $-2.33$.
\end{solution}

\begin{solution}{exo:b3:water-rock-interaction:12}
$1 - 10^{\mathrm{SI}}$: 0.99, 0.90 dan 0.21. Hukum pangkat memberikan
$(10^4/6)^{-0.61} = 0.011$: laju sekitar 90 kali lebih rendah. Mineral di lapangan
sudah tua, permukaan reaktifnya habis, dan direndam air pori yang lambat
serta mendekati kesetimbangan; kedua faktor ini saling melipatgandakan, sehingga laju di lapangan turun beberapa orde besaran
di bawah laju laboratorium.
\end{solution}

\begin{solution}{pb:b3:water-rock-interaction:1}
\textbf{1.} \ce{Ca} 2.795, \ce{Mg} 0.165, \ce{Na} 0.305, \ce{K} 0.031,
\ce{HCO3} 5.572, \ce{Cl} 0.310, \ce{SO4} 0.146, \ce{NO3} \qty{0.145}{mmol/L}.

\textbf{2.} Kation \qty{6.25}{meq/L}, anion \qty{6.32}{meq/L}: galat
\qty{-0.5}{\percent}, dapat diterima.

\textbf{3.} $I = \tfrac12(4\times2.795 + 4\times0.165 + 0.305 + 0.031 + 5.572
+ 0.310 + 4\times0.146 + 0.145)\times10^{-3} = 0.0094$.

\textbf{4.} $\sqrt I = 0.0969$: $\gamma_{\ce{Ca^2+}} = 0.683$,
$\gamma_{\ce{HCO3-}} = 0.904$, $\gamma_{\ce{CO3^2-}} = 0.676$.

\textbf{5.} Ya: $I$ hanya sepersepuluh dari batas sekitar 0.1.

\textbf{6.} $a_{\ce{HCO3-}} = 5.572\times10^{-3}\times0.904 = 5.04\times
10^{-3}$ ($\log = -2.298$).

\textbf{7.} $a_{\ce{CO2}} = 10^{-7.20}\times5.04\times10^{-3}/10^{-6.46} =
9.17\times10^{-4}$ ($\log = -3.038$).

\textbf{8.} $\log P_{\ce{CO2}} = -3.038 + 1.27 = -1.77$, \qty{0.017}{atm}:
$10^{1.60} = 40$ kali udara, berada dalam kisaran air tanah. Karbon
dioksida berasal dari udara tanah di daerah imbuhan, dari akar dan dekomposisi.

\textbf{9.} $a_{\ce{CO3^2-}} = 10^{-10.49}\times5.04\times10^{-3}/10^{-7.20}
= 2.58\times10^{-6}$, $[\ce{CO3^2-}] = \qty{3.8e-6}{mol/L}$. Total
\qty{6.49}{mmol/L}: \ce{CO2} \qty{14.1}{\percent}, \ce{HCO3-}
\qty{85.8}{\percent}, \ce{CO3^2-} \qty{0.06}{\percent}.

\textbf{10.} Alkalinitas $5.572 + 2\times0.0038 = \qty{5.580}{meq/L}$:
\qty{99.9}{\percent} bikarbonat.

\textbf{11.} $a_{\ce{Ca^2+}} = 2.795\times10^{-3}\times0.683 = 1.909\times
10^{-3}$; $\log\mathrm{IAP} = \log(1.909\times10^{-3}\times2.58\times10^{-6})
= -8.307$; $\mathrm{SI} = -8.307 + 8.35 = +0.04$.

\textbf{12.} Dalam $\pm0.1$, ketidakpastian indeks: air berada dalam
kesetimbangan dengan kalsit. Air telah lama berada di batuan: mata air yang
lambat dan menyebar.

\textbf{13.} $m = [10^{-5.59}\times0.0171/(4\times0.683\times0.904^2)]^{1/3}
= \qty{2.70}{mmol/L}$, \qty{108}{mg/L}, dibandingkan \qty{2.79}{mmol/L} terukur:
\qty{3}{\percent} lebih tinggi, sesuai dengan indeks $+0.04$.

\textbf{14.} $a_{\ce{CO2}} = 10^{-1.27 - 3.37} = 2.29\times10^{-5}$;
$a_{\ce{H+}} = 10^{-6.46}\times2.29\times10^{-5}/5.04\times10^{-3}$, pH 8.80;
$a_{\ce{CO3^2-}} = 1.03\times10^{-4}$; $\mathrm{SI} = \log(1.909\times
10^{-3}\times1.03\times10^{-4}) + 8.35 = +1.65$ (1.61 dengan
neraca alkalinitas penuh pada gambar).

\textbf{15.} Iterasi $m = [10^{-5.59 - 3.37}/(4\gamma_{\ce{Ca}}
\gamma_{\ce{HCO3}}^2)]^{1/3}$ dengan $I = 3m$: $m = \qty{0.72}{mmol/L}$
(\qty{29}{mg/L}), $I = 0.0022$. Kalsit mengendap: $2.795 - 0.718 =
\qty{2.08}{mmol/L}$, sekitar \qty{210}{mg} per liter.

\textbf{16.} $0.208\times2.0 = \qty{0.42}{g/s}$, \qty{13}{t} per tahun,
$13.1/2.71 = \qty{4.8}{m^3}$; di atas \qty{400}{m^2}, \qty{12}{mm} per tahun,
\qty{0.03}{mm} per hari: jauh di bawah tufa mata air dingin yang disebutkan dalam bab.

\textbf{17.} $X = \log(1.2\times10^{-3}/39.10) + 5.8 = 1.29$; $Y =
\log(9.0\times10^{-3}/60.08) = -3.82$.

\textbf{18.} $X < 1.55$ dan $Y > -4.39$: kaolinit. Kuarsa: $\mathrm{SI} =
-3.82 + 3.98 = +0.16$.

\textbf{19.} Di bawah $\log a_{\ce{H4SiO4}} = -4.39$, terdapat \qty{2.4}{mg/L}
silika.

\textbf{20.} $X = 1.55$ memerlukan pH $= 1.55 + 4.51 = 6.06$.

\textbf{21.} Asam $-10.06 - 0.500\times5.8 = -12.96$; netral $-12.41$; basa
$-21.20 + 0.823\times5.8 = -16.43$. Jumlah: $\log r = -12.30$, suku netral
dominan.

\textbf{22.} Terbatas-pasokan: tanah tebal dan lereng landai berarti erosi lambat;
feldspar segar yang mencapai front pelapukan terlindi penuh, dan
fluks mengikuti laju erosi, bukan iklim.

\textbf{23.} $\exp[(74/0.008314)(1/283.15 - 1/298.15)] = 4.9$.

\textbf{24.} \omterm{def:b3:water-rock-interaction:si}{Indeks kejenuhan} kalsit pada air mata air di keluarannya
adalah $\mathrm{SI} = \textbf{+0.04}$.
\end{solution}
